\documentclass[12pt]{article}
\usepackage[utf8]{inputenc}
\usepackage[T1]{fontenc}
\usepackage{hyperref}
\title{Taiwan Sign Language Language Resources for K-12 Curriculum Development}
\author{Compiled by Claude Sonnet 5 (claude-sonnet-5) for Andr\'as Kornai}
\date{}
\begin{document}
\maketitle

\section{Introduction}
Taiwan Sign Language (TSL; ISO 639-3 tss; Glottocode taiw1241) is the language of Taiwan's deaf community. It was included in the Development of National Languages Act in 2019. The Ministry of Education began incorporating TSL courses into school curricula in 2022, according to a January 2026 Taiwan Today report seen in search results. The ministry also launched 34 distance-learning TSL classes taught by deaf instructors starting in October 2025. Deaf schools have used Total Communication since the early 1970s.

\section{Primary Educational Resources}
\begin{itemize}
\item ``Literacy and deaf students in Taiwan: Issues, practices and directions for future research, Part I'' (NTNU scholar database), on Chinese literacy teaching in deaf education: \url{https://scholar.lib.ntnu.edu.tw/en/publications/literacy-and-deaf-students-in-taiwan-issues-practices-and-directi-4/}
\item Part II of the same study: \url{https://scholar.lib.ntnu.edu.tw/en/publications/literacy-and-deaf-students-in-taiwan-issues-practices-and-directi-3/}
\end{itemize}
No TSL textbook series for the K-12 curriculum was found. Earlier TSL textbooks (Smith and Ting 1979, 1984) were sponsored by the National Association of the Deaf of the Republic of China.

\section{Linguistic and Reference Materials}
\begin{itemize}
\item Gallaudet Visual Language and Visual Learning (VL2) center paper, which may concern TSL (the search snippet did not confirm the subject): \url{https://vl2.gallaudet.edu/center-papers/233}
\end{itemize}
A 2005 article by Wayne Smith listed a comprehensive TSL dictionary as future work. No comprehensive dictionary was confirmed.

\section{Religious and Community Resources}
No community or religious resources were identified in the searches.

\section{Digital and Audio Resources}
The searches found no TSL dictionary app or video corpus.

\section{Curriculum Development Context}
There is official support from the 2019 law and the ministry's 2022 curriculum inclusion, but deaf-taught TSL classes only began as remote classes in October 2025 because of teacher shortages. Next steps include locating the MOE curriculum guidelines and the existing TSL dictionaries, which the searches did not find.
\end{document}